\section{Summary}
We have described a method that resolves both the theoretical
and numerical challenges faced when calculating moments
of parton distribution functions from lattice QCD.
A lattice regulator breaks the
Euclidean rotational symmetry into the hypercubic group.
As a result, the renormalization of twist-2 operators
on the lattice is a substantial hurdle to determine
moments of parton distribution functions
$\left\langle x^{n-1} \right\rangle$.
While for $n>4$, power divergences in the
lattice spacing are unavoidable, for $n=3,4$
the cancellation of power divergences is still possible
by choosing suitable irreducible representations
of the hypercubic group. The price to pay is that
the calculation of the hadronic matrix elements
necessitates the use
of a spatial external momentum, worsening the signal-to-noise
ratio of the correlation functions.
At finite flow time $t$, the correlation functions
of twist-2 operators
are finite once the bare parameters of the theory and the flowed fermion
fields are renormalized.
After performing the continuum limit,
the matching with the physical matrix element can be done
at short flow time using the continuum symmetries of the theory.
Using traceless and symmetrized operators guarantees that the
matching is multiplicative, and we have calculated, at
1-loop in perturbation theory, the matching coefficients.

It would be also very interesting
to investigate how the approach described in this work
can complement recent efforts aimed at the direct
determination of PDFs from lattice QCD correlators.

While, the method proposed in this paper is applicable,
at least in principle, to calculate any moments
of parton distribution functions,
in practice lattice QCD calculations of hadronic matrix elements of
the flowed twist-2 operators in Eq.~\eqref{eq:flowed_t2} are subject to
several systematic uncertainties.

In the Supplemental Material~\cite{SupplMat} there is a discussion of the
systematics, but only a thorough numerical investigation
provides the final response about the
range of applicability, and whether
this result opens the way to determine moments of any order
of the parton distribution functions.
